\section{Related Work}
\label{sec:related}

\para{The Linear Representation Hypothesis.}
The idea that concepts correspond to directions in activation space has been formalized by \citet{park2023linear}, who unify three interpretations---subspace, measurement, and intervention---under a causal inner product framework. Testing 27 concepts on LLaMA-2-7B, they find linear structure for 26, with the sole exception being the relational concept \texttt{thing$\rightarrow$part}. \citet{jiang2024origins} provide a theoretical account: linearity arises from the interaction of next-token prediction with softmax and cross-entropy, but this analysis assumes binary concepts with clear contrasts. \citet{park2024geometry} extend the \lrh to categorical concepts (polytopes) and hierarchical concepts (orthogonal subspaces), showing that even within the linear framework, complex concept types require richer geometric structures. Our work tests whether the \lrh holds for emotional styles and, crucially, whether linear detectability implies linear controllability.

\para{Steering vectors and activation engineering.}
\citet{turner2023steering} introduce Activation Addition (\actadd), computing steering vectors from contrastive prompt pairs and adding them to the residual stream. They demonstrate effective steering for sentiment and topics but observe failures for compositional concepts and non-monotonic behavior when using partial dimensions. \citet{rimsky2024steering} scale this approach via Contrastive Activation Addition (\caa), averaging activation differences across many examples and testing on behavioral dimensions (sycophancy, corrigibility) in Llama~2 Chat. \citet{konen2024style} apply both activation-based and training-based style vectors to emotions and writing styles in Alpaca-7B, finding that disgust and surprise resist steering while joy and anger succeed. Unlike these works, we systematically compare steering effectiveness against linear separability and prompting success across the same style taxonomy.

\para{Non-linear representations.}
\citet{csordas2024nonlinear} discover ``onion representations'' in small GRUs, where tokens at different positions share the same direction but differ in magnitude, requiring nonlinear decoding. Linear probing and distributed alignment search completely fail on these representations while nonlinear interventions succeed. \citet{elhage2022superposition} show that neural networks store more features than dimensions via superposition, organizing them into geometric polytope structures. This implies that features sharing dimensions cannot be cleanly separated by linear methods. Whether similar non-linear encodings exist for style in large transformers remains unexplored---our probing experiments directly test this.

\para{Benchmarking steering methods.}
\citet{wu2025axbench} introduce AxBench, the first large-scale benchmark comparing steering methods across 500 concepts in Gemma-2-2B/9B. Their key finding is that prompting outperforms all representation-based steering methods, including sparse autoencoders and diff-in-means vectors. This ``prompting dominance'' is consistent with our results but does not explain \emph{why} prompting succeeds where linear methods fail. Our work provides a mechanistic perspective: the detection--control gap suggests that prompting accesses style through computational pathways that bypass the limitations of single-direction activation addition. \citet{han2023word} find that style control through output embedding transformations requires a full $d \times d$ matrix rather than a single direction, further suggesting that style involves higher-dimensional structure than a single linear direction can capture.

\para{Style transfer and controlled generation.}
The broader literature on style transfer and controlled generation provides context for our style taxonomy. \citet{demszky2020goemotions} introduce GoEmotions, the dataset we use, with 54K Reddit comments annotated for 27 emotion categories. \citet{chalnev2024improving} use sparse autoencoders to measure causal effects of steering vectors and find that naive linear steering produces unpredictable side effects, consistent with our observation that the relationship between activation-space metrics and behavioral effects is nonlinear.
